% Bản nháp chương 8 theo dàn bài mới, viết với bộ mẫu ở style/mau/.
% ===========================================================================
\chapter{Kết luận bền tới đâu}
\label{sec:lambda}

Mọi con số trong bốn chương trước đo ở đúng một giá trị $\lambda$, tức ở đúng một
số điều kiện $\kappa = \num{268.31}$. Vì $\kappa$ có mặt trong cả hai cận ở
chương~\ref{sec:theory} và trong cả ba cơ chế ở chương~\ref{sec:mechanisms}, câu
hỏi tự nhiên là các kết luận ấy còn đúng tới đâu khi bài toán dễ đi hoặc khó đi.
Chương này đổi $\lambda$ để đổi $\kappa$ và đo lại.

\section{Hệ số hiệu chỉnh điều khiển số điều kiện}
\label{sec:lambda-control}

\resultfig{lambda_effect_iter}
  {Tốc độ hội tụ ứng với ba giá trị $\lambda$, theo số vòng lặp. Cả ba lần chạy
   dùng gradient descent với $t = 1/L$.}
  {fig:lambda-iter}
\resultfig{lambda_effect_time}
  {Cùng ba lần chạy ấy theo thời gian chạy.}
  {fig:lambda-time}

\begin{table}[tp]
  \centering
  \caption{Hệ số hiệu chỉnh vừa quyết định số điều kiện, vừa quyết định chất
    lượng dự báo. Cả ba lần chạy dùng gradient descent với $t = 1/L$, nên cột số
    vòng lặp đặt cạnh được với dòng $t = 1/L$ của bảng~\ref{tab:gd-fixed}. Cột
    cuối tính theo \eqref{eq:rmse-to-price} trên một chiếc máy giá trung vị.}
  \label{tab:lambda}
  \begin{tabular}{rrrrrrr}
    \toprule
    $\lambda$ & $\kappa$ & \multicolumn{2}{c}{Vòng lặp đạt}
      & Tổng (s) & RMSE test & Sai lệch giá \\
    \cmidrule(lr){3-4}
    & & $10^{-6}$ & $10^{-12}$ & & & \\
    \midrule
    \num{0.001} & \num{2674}  & 40 & 170 & \num{10.58} & \num{0.20580} & \num{-4.8}  \\
    \num{0.01}  & \num{268.3} & 40 & 150 & \num{8.42}  & \num{0.20601} & \num{0}     \\
    \num{0.1}   & \num{27.73} & 30 & 90  & \num{5.46}  & \num{0.21916} & \num{302}   \\
    \bottomrule
  \end{tabular}
\end{table}

Ba dòng của bảng~\ref{tab:lambda} cho $\kappa$ bằng \num{2674}, \num{268.3} và
\num{27.73} ứng với $\lambda$ bằng \num{0.001}, \num{0.01} và \num{0.1}, tức tăng
$\lambda$ lên 10 lần thì giảm $\kappa$ đúng 10 lần. Quan hệ ấy đến thẳng từ
\eqref{eq:mu-equals-lambda}: trên bài toán này $\mu = \lambda$ nên $\kappa$ là
$L/\lambda$, và $L$ không đổi khi $\lambda$ đổi vì $\lambda$ chỉ cộng thêm một
lượng nhỏ so với trị riêng lớn nhất. Hiệu chỉnh Ridge do đó có hai vai trò tách
biệt, một vai trò thống kê là chống quá khớp và một vai trò tối ưu hóa là làm bài
toán dễ giải hơn.

\section{Số điều kiện là chỉ báo kém cho chi phí}
\label{sec:lambda-early}

Hai vai trò ấy không cùng chiều, và mức ảnh hưởng của chúng lệch nhau rất xa. Từ
$\lambda = \num{0.001}$ lên $\num{0.1}$, số điều kiện giảm 96 lần nhưng số vòng
lặp cần để đạt $10^{-6}$ chỉ giảm từ 40 xuống 30, tức 25\%
(hình~\ref{fig:lambda-iter}). Nếu $\kappa$ chi phối chi phí như \eqref{eq:gd-rate}
gợi ý thì mức giảm phải gần với 96 lần chứ không phải một phần tư.

Nguyên nhân đã có sẵn ở chương~\ref{sec:theory}. Ngưỡng $10^{-6}$ rơi vào giai
đoạn đầu, nơi sai số do 269 hướng cong chi phối và chúng tắt theo $\abs{1 - tL}$
chứ không theo $(\kappa-1)/(\kappa+1)$. Đổi $\lambda$ không đụng tới các hướng ấy
vì nó chỉ nâng đầu nhỏ của phổ, nên nó gần như không đổi được chi phí ở mức chính
xác mà bài toán định giá cần.

\section{Ở độ chính xác cao thì số điều kiện mới lên tiếng}
\label{sec:lambda-tail}

Cùng ba lần chạy ấy nhưng đo tới $10^{-12}$ thì bức tranh đổi. Số vòng lặp đi từ
170 xuống 90, tức giảm 47\% thay vì 25\%, và tổng thời gian tới lúc dừng hẳn đi
từ \SI{10.58}{\second} còn \SI{5.46}{\second} (hình~\ref{fig:lambda-time}). Độ
nhạy với $\kappa$ tăng lên khi ngưỡng chặt lại, đúng như bức tranh hai giai đoạn
ở mục~\ref{sec:theory-early} dự đoán.

Nhưng ngay cả ở $10^{-12}$, giảm $\kappa$ đi 96 lần cũng mới cắt được nửa số vòng
lặp. Muốn thấy đúng tỉ lệ mà \eqref{eq:gd-rate} mô tả thì phải chạy tới chỗ chỉ
còn tám hướng phẳng chi phối, và trên bài toán này chỗ đó nằm sát giới hạn số học
của kiểu dữ liệu, tức nằm ngoài vùng có ý nghĩa thực tế. Nói cách khác, $\kappa$
là đại lượng đúng để dự đoán hành vi tiệm cận nhưng là chỉ báo kém cho chi phí ở
mọi mức chính xác mà người ta thực sự dùng.

\section{Cái giá về phía bài toán định giá}
\label{sec:lambda-price}

Cột cuối bảng~\ref{tab:lambda} cho thấy vế còn lại của đánh đổi. Từ
$\lambda = \num{0.001}$ lên $\num{0.1}$, RMSE trên tập kiểm tra xấu đi 6,5\%, và
quy theo \eqref{eq:rmse-to-price} thì sai lệch giá tăng \num{307} đơn vị tiền
trên một chiếc máy giá trung vị.

Con số ấy đáng đặt cạnh mọi con số khác trong báo cáo. Chênh lệch giữa cấu hình
tối ưu hóa tốt nhất và tệ nhất mà nhóm tự cài là 87,7 đơn vị tiền, còn giữa nhóm
và cài đặt mặc định tệ nhất của thư viện là \num{137} đơn vị
(bảng~\ref{tab:sklearn}). Một quyết định về hệ số hiệu chỉnh vì thế đắt hơn mọi
quyết định về thuật toán cộng lại, dù nó chỉ là một dòng trong phần chuẩn bị dữ
liệu.

Đổi lại, khoản thu được về phía tối ưu hóa là 10 vòng lặp ở ngưỡng $10^{-6}$, tức
khoảng \SI{0.3}{\second}. Đánh đổi ấy không có lợi trên bài toán này, và đó là lý
do quy tắc một sai số chuẩn ở mục~\ref{sec:choose-lambda} dừng ở
$\lambda = \num{0.01}$ chứ không đi xa hơn về phía hiệu chỉnh mạnh.

\section{Kết luận nào bền, và điều chưa kiểm được}
\label{sec:lambda-scope}

Ba kết luận của báo cáo giữ nguyên qua hai bậc thay đổi của $\kappa$. Ngưỡng
$2/L$ không phụ thuộc $\lambda$ vì $L$ không đổi. Bức tranh hai giai đoạn ở
chương~\ref{sec:theory} giữ nguyên và còn được củng cố, vì độ nhạy với $\kappa$
tăng đúng theo hướng nó dự đoán khi ngưỡng chặt lại. Và ngưỡng
$\eps_{\text{app}}$ của chương~\ref{sec:enough} vẫn nằm trong giai đoạn đầu ở cả
ba giá trị $\lambda$, nên kết luận rằng chọn thuật toán không đổi được chất lượng
định giá cũng không đổi.

Phần chưa kiểm được cần nói rõ. Thí nghiệm này chạy riêng gradient descent ở ba
giá trị $\lambda$, nên nó trả lời câu $\kappa$ dự đoán chi phí tốt tới đâu, chứ
chưa trả lời câu thứ hạng giữa các thuật toán ở bảng~\ref{tab:summary} có đổi khi
$\kappa$ đổi hay không. Muốn trả lời câu sau phải chạy lại toàn bộ lưới tham số ở
cả ba giá trị $\lambda$, tức gấp ba khối lượng tính toán của
chương~\ref{sec:mechanisms}. Nhóm dự đoán thứ hạng giữ nguyên, vì cả ba cơ chế
chi phối nó đều không phụ thuộc $\lambda$, nhưng đây là dự đoán chứ chưa phải kết
quả đo.
